\documentclass[aspectratio=169,11pt]{beamer}

\usepackage{fontspec}
\usepackage{graphicx}
\usepackage{xcolor}
\usepackage{hyperref}
\usepackage{kotex}

\defaultfontfeatures{Ligatures=TeX,Scale=MatchLowercase}
\IfFontExistsTF{Noto Sans CJK KR}{
  \setmainfont{Noto Sans CJK KR}
  \setsansfont{Noto Sans CJK KR}
}{
  \IfFontExistsTF{Apple SD Gothic Neo}{
    \setmainfont{Apple SD Gothic Neo}
    \setsansfont{Apple SD Gothic Neo}
  }{
    \setmainfont{NanumGothic}
    \setsansfont{NanumGothic}
  }
}
\IfFontExistsTF{DejaVu Sans Mono}{
  \setmonofont{DejaVu Sans Mono}
}{
  \IfFontExistsTF{Menlo}{
    \setmonofont{Menlo}
  }{
    \setmonofont{Liberation Mono}
  }
}

\setbeamertemplate{navigation symbols}{}
\setbeamertemplate{footline}[frame number]
\setbeamertemplate{itemize items}[circle]
\setbeamertemplate{blocks}[rounded][shadow=false]
\setbeamercolor{title}{fg=black}
\setbeamercolor{frametitle}{fg=black}
\setbeamercolor{block title}{fg=black,bg=blue!8}
\setbeamercolor{block body}{fg=black,bg=gray!6}

\title{UUV MuJoCo v2.2\\Fluid Proxy Update}
\author{}
\date{2026-04-07}

\begin{document}

\begin{frame}[t]{Ellipsoid Fluid Proxy 설계}
  \small
  \begin{columns}[T,onlytextwidth]
    \begin{column}{0.59\textwidth}
      \includegraphics[width=\linewidth]{figures_v22_latest/uuv_v22_latest_fluid_proxy_overview.png}
    \end{column}
    \begin{column}{0.39\textwidth}
      \begin{block}{핵심 설명}
        \begin{itemize}
          \item MuJoCo built-in fluid는 실제 메쉬 전체가 아니라 단순화한 proxy 형상으로 유체력을 계산한다.
          \item current는 중앙 enclosure, 좌우 lower body, 상하 plate proxy로 구성했다.
          \item \texttt{fluidcoef = [blunt, slender, angular, kutta, magnus]} 이며, center는 부피와 회전 저항을, lower body는 전진 및 yaw 저항을 맡는다.
        \end{itemize}
      \end{block}
      \vspace{0.3em}
      \begin{block}{핵심 문장}
        중앙 body는 부피와 회전 저항을, 좌우 body는 전진 및 yaw 저항을, 상하 plate는 heave 저항 보강을 담당한다.
      \end{block}
    \end{column}
  \end{columns}
\end{frame}

\begin{frame}[t]{왜 Plate Proxy를 추가했는가}
  \small
  \begin{columns}[T,onlytextwidth]
    \begin{column}{0.60\textwidth}
      \includegraphics[width=\linewidth]{figures_v22_latest/uuv_v22_latest_plate_proxy_compare.png}
    \end{column}
    \begin{column}{0.38\textwidth}
      \begin{block}{수정 이유}
        \begin{itemize}
          \item 기존 lower body는 ellipsoid만 사용해서 상하가 너무 둥글게 표현되었다.
          \item 그래서 실제보다 heave 방향 투영면적이 작게 잡혀 z축 저항이 약하게 나올 수 있었다.
          \item 이를 보완하기 위해 좌우 lower body 위와 아래에 얇은 plate proxy를 추가했고, plate는 더 작은 계수로 시작해 heave 감쇠만 보강하도록 두었다.
        \end{itemize}
      \end{block}
      \vspace{0.3em}
      \begin{block}{핵심 문장}
        plate proxy는 평평한 상하 면적을 보강해, 기존 ellipsoid-only 구성에서 부족했던 z축 유체저항을 보완한다.
      \end{block}
    \end{column}
  \end{columns}
\end{frame}

\end{document}
